\begin{center}
    \footnotesize
    \renewcommand{\arraystretch}{1.1}
    \rowcolors{2}{gray!8}{white}

    \begin{longtable}{l r r}
        \toprule
        \textbf{Subspecies} & $\boldsymbol{\mu\,(10^{-8}\text{/site/gen})}$ & $\boldsymbol{N_e\,(10^{4})}$ \\
        \midrule
        \endfirsthead
        \multicolumn{3}{c}{\tablename\ \thetable{} -- continued from previous page} \\
        \toprule
        \textbf{Subspecies} & $\boldsymbol{\mu\,(10^{-8}\text{/site/gen})}$ & $\boldsymbol{N_e\,(10^{4})}$ \\
        \midrule
        \endhead
        \bottomrule
        \addlinespace[2ex]
        \rowcolor{white}\caption{Per-subspecies mutation rate and inferred effective population size. Mutation rates $\mu$ (per site per generation) are taken from \citet{kuderna2023} at the species level and applied to all subspecies of that species; for the four subspecies absent from \citet{kuderna2023} (\textit{T.\ poliocephalus}, \textit{M.\ hecki}, \textit{M.\ leucogenys}, \textit{M.\ assamensis}), $\mu$ is the value reported for the reference species onto which they were mapped. Effective population sizes $N_e$ are estimated as $\hat\theta / (4\mu)$ using Watterson's estimator on neutral (4-fold) sites (see Methods). Subspecies are ordered by $N_e$.}\label{tab:species_summary}
        \endlastfoot
        \textit{Gorilla beringei} graueri & 1.20 & 1.49 \\
        \textit{Pan paniscus} & 1.17 & 1.60 \\
        \textit{Homo sapiens} & 1.20 & 1.64 \\
        \textit{Pan troglodytes} verus & 1.31 & 1.73 \\
        \textit{Trachypithecus poliocephalus} & 0.72 & 2.04 \\
        \textit{Gorilla beringei} beringei & 1.20 & 2.12 \\
        \textit{Rhinopithecus bieti} & 0.52 & 2.13 \\
        \textit{Pan troglodytes} ellioti & 1.31 & 2.69 \\
        \textit{Trachypithecus francoisi} & 0.72 & 2.93 \\
        \textit{Pongo abelii} & 1.47 & 3.13 \\
        \textit{Pan troglodytes} schweinfurthii & 1.31 & 3.45 \\
        \textit{Pongo pygmaeus} & 1.28 & 3.51 \\
        \textit{Gorilla gorilla} gorilla & 1.33 & 3.60 \\
        \textit{Macaca silenus} & 0.55 & 3.90 \\
        \textit{Pan troglodytes} troglodytes & 1.31 & 4.35 \\
        \textit{Rhinopithecus roxellana} & 0.49 & 4.92 \\
        \textit{Theropithecus gelada} & 0.44 & 5.87 \\
        \textit{Papio hamadryas} & 0.61 & 6.64 \\
        \textit{Papio ursinus} & 0.69 & 7.40 \\
        \textit{Macaca nigra} & 0.49 & 7.72 \\
        \textit{Macaca cyclopis} & 0.44 & 8.28 \\
        \textit{Macaca radiata} & 0.44 & 8.76 \\
        \textit{Papio kindae} & 0.72 & 9.17 \\
        \textit{Macaca fuscata} & 0.43 & 9.36 \\
        \textit{Papio cynocephalus} anubishybrid & 0.59 & 10.14 \\
        \textit{Macaca maura} & 0.54 & 10.37 \\
        \textit{Macaca arctoides} & 0.48 & 11.27 \\
        \textit{Semnopithecus hypoleucos} & 0.49 & 11.60 \\
        \textit{Papio cynocephalus} ssp & 0.59 & 11.86 \\
        \textit{Erythrocebus patas} & 0.49 & 13.25 \\
        \textit{Macaca leonina} & 0.51 & 13.67 \\
        \textit{Papio anubis} & 0.39 & 13.86 \\
        \textit{Macaca mulatta} & 0.56 & 15.93 \\
        \textit{Macaca hecki} & 0.37 & 16.04 \\
        \textit{Macaca tonkeana} & 0.53 & 16.05 \\
        \textit{Macaca leucogenys} & 0.39 & 17.16 \\
        \textit{Macaca assamensis} & 0.39 & 17.84 \\
        \textit{Macaca nemestrina} & 0.37 & 28.17 \\
    \end{longtable}
\end{center}
